\documentclass[10pt,a4paper]{amsart}
\usepackage[margin=19mm]{geometry}
\usepackage{amsmath,amssymb,mathtools}
\usepackage[scr=rsfs,cal=euler]{mathalfa}
\usepackage{xfrac}
\usepackage[unicode,hidelinks,bookmarks=false]{hyperref}
\newcommand{\ie}{i.e.\ }
\newcommand{\cf}{cf.\ }
\newcommand{\resp}{respectively\ }
\newcommand{\Subsection}{\subsection}
\providecommand{\nobreakdash}{\nobreak\hbox{-}}
\setlength{\emergencystretch}{2em}
\multlinegap0pt
\usepackage[normalem]{ulem}
\usepackage{amssymb}
\usepackage{mathtools}
\usepackage{tikz}
\newtheorem{proposition}{Proposition}
\newcommand{\Ch}{\mathrm{Ch}}
\newcommand{\R}{\mathbb{R}}
\newcommand{\sfd}{\mathsf d}
\newcommand{\vol}{\mathrm{vol}}
\title{The infinitesimal structure of manifolds with non-continuous Riemannian metrics}
\author{Vanessa Ryborz}
\date{Source: arXiv:2507.14726v3 (CC BY 4.0)}
\begin{document}
\maketitle

\noindent\emph{Source and licence.} The infinitesimal structure of manifolds with non-continuous Riemannian metrics. Version arXiv:2507.14726v3 (CC BY 4.0), distributed by arXiv under the Creative Commons Attribution 4.0 International licence, http://creativecommons.org/licenses/by/4.0/, as recorded in the arXiv licence metadata for this exact version and shown on the abstract page https://arxiv.org/abs/2507.14726v3. Full licence notice: \texttt{licenses/arxiv-2507.14726-CC-BY-4.0.txt}.\par\smallskip

\noindent\textit{Editorial scope.} Counted unit: the article's proposition first\_non\_infinitesimally\_hilb (source0.tex lines 918--920). The article's complete original proof of it is delivered with it (source0.tex lines 921--950, one \texttt{proof} environment of 30 lines). No mathematics is written, repaired or re-derived: the statement and the proof are the article's own lines, in the article's own notation, and no retained line is substituted or re-flowed.  No further source-local context is needed: every cross-reference of every retained line resolves inside the two delivered spans.  Nothing was omitted from any retained span, so the package carries no omission record.  No retained line contains a citation, so the wrapper carries no bibliography and no bibliography entry.  No retained line contains a figure environment, an image-inclusion command or drawing code, so no image is delivered and the package declares no asset dependency.  Editorial formatting only, in addition: an AMS \texttt{amsart} preamble that restates the article's own declarations and macros used by the retained lines, namely the environments \texttt{proposition} on separate counters, together with the standard packages those lines need.  The retained environments print in this document as Proposition 1 in order of appearance, whereas the article prints the same items with the numbering of its own full document.  The complete native source of the article as distributed in the arXiv e-print for this exact version is bundled unchanged as \texttt{sources/181874/source0.tex}.
\begin{proposition}\label{first_non_infinitesimally_hilb}
    The metric measure space $(\R^d, \sfd_g, \vol_g)$ is not infinitesimally Hilbertian. 
\end{proposition}

\begin{proof}
    Pick a function $\phi \in C_c^\infty(\R^d)$ such that $\phi \equiv 1$ on $[-1, 2]^d$. Consider the functions $\Tilde{f}_i := \phi f_i$ for $1 \leq i \leq 4$. Then by the previous observations, we have that 
    \begin{align*}
        &|D\Tilde{f}_i|_w = G_i \ \mathrm{in}\ (-1, 2)^d\\
        &|D\Tilde{f}_i|_w = \frac{1}{\sqrt{2}}|\nabla_{euc} \Tilde{f}_i|_{euc} \ \mathrm{in}\ ((-1,2)^d)^c.
    \end{align*}
    More precisely, this gives
    \begin{align}
        |D\Tilde{f}_i|_w= |\nabla_g \Tilde{f}_i|_g \ \mathrm{a.e.\ in} \ (O\times(0,1)^{d-1}) \cup ((0,1)^{d})^c,
    \end{align}
    and
    \begin{align*}
        &|D\Tilde{f}_1|_w= \frac{1}{\sqrt{2}}, \ \mathrm{a.e.\ in} \ (0,1)^{d}\setminus O\times(0,1)^{d-1}, \\
        &|D\Tilde{f}_2|_w=|D\Tilde{f}_3|_w=|D\Tilde{f}_4|_w=1 \ \mathrm{a.e.\ in} \ (0,1)^{d}\setminus (O\times(0,1)^{d-1}).
    \end{align*}
    Then it immediately follows that 
    \begin{align}
    |D\Tilde{f}_3|_w^2 + |D\Tilde{f}_4|_w^2 = 2(|D\Tilde{f}_1|_w^2 + |D\Tilde{f}_2|_w^2)   \ \mathrm{a.e.\ in} \ (O\times(0,1)^{d-1}) \cup ((0,1)^{d})^c,
    \end{align}
    but 
    \begin{align}
        |D\Tilde{f}_3|_w^2 + |D\Tilde{f}_4|_w^2 - 2(|D\Tilde{f}_1|_w^2 + |D\Tilde{f}_2|_w^2) = -1 \ \mathrm{a.e.\ in} \ (0,1)^{d}\setminus (O\times(0,1)^{d-1}).
    \end{align}
    Moreover, we have that $\Tilde{f}_3 = \Tilde{f}_1 + \Tilde{f}_2$ and $\Tilde{f}_4 = \Tilde{f}_1 - \Tilde{f}_2$ and all four functions have finite Cheeger energy.
    But as $\mathcal{L}^d((0,1)^{d}\setminus O\times(0,1)^{d-1})\geq 1-2\kappa >0$ and the fact that $\vol_g$ and $\mathcal{L}^d$ are $2^{\frac{d}{2}}$-equivalent, we get that 
    \begin{align}
        \Ch(\Tilde{f}_3) + \Ch(\Tilde{f}_4) \neq 2(\Ch(\Tilde{f}_1) + \Ch(\Tilde{f}_2)),
    \end{align}
    which proves that the metric measure space $(\R^d, \sfd_g,  \vol_g)$ is not infinitesimally Hilbertian.
\end{proof}

\end{document}
